\cotefile{126}
\scanpage{2}
% archive sheet number: 1

% slanted
\margin{Compléter, cf remarque 2 \uline{plus bas}}

\uline{Théorème} 1 \quad $f:X^{\unsure{\bullet|\circ|{}}}\longrightarrow S$ un morphisme \uline{\unsure{fid.} plat quasi-cpt}, \inserted{\uline{quasi-séparé}, \deleted{S} \uline{affine}} \inserted{\uline{couvrant quasi-compact quasi-séparé plat}}, $F$ faisceau quasi-cohérent sur $X$. Alors

\[
\unsure{H^i|\check H^i}(X^{\bullet}/S,F)\simeq
\begin{cases}0&\text{si }i\ne0\\ H^0(S,F)&\text{si }i=0\end{cases}
\]

On peut l'énoncer sous la forme suivante

\uline{Corollaire} 1 \quad \unsure{Hypothèse|Hypothèses} \unsure{précédente|précédentes}, sans hyp. affine sur $X$ cependant. Alors $\underline H^i(X^{\bullet}/S,F)$ (faisceaux \inserted{qu. coh.} sur $X$) \deleted{sont} \deleted{\unsure{isomorphes}} sont \unsure{nul|nuls} si $i\ne0$, et $F\longrightarrow\underline H^0(X\unsure{'|\bullet|^{\cdot}}/X,F)$ est un isomorphisme.

Or la formation \deleted{des} $\underline H^i(X^{\bullet}/S,F)$ est compatible avec extension \uline{plate} de \uline{la base}. [car les $X^{i}/S$ sont quasi-compacts quasi-séparés], et il suffit de prouver que les assertions dites sont vraies après changement de base plat couvrant $S_\lambda\longrightarrow S$, \inserted{les $S_\lambda$ affines.} Or on peut en trouver qui sont tels que les $X_\lambda/S_\lambda$ aient des sections, d'où le résultat, car pour $S$ affine et $X/S$ ayant une section, c'est immédiat.

% margin note written sideways
\margin{\uline{Corollaire} 2 Soit $X\longrightarrow S$ \uline{couvrant} \deleted{\unsure{q-comp.}} $H^0(S,F)\simeq \unsure{\check H^0|H^0}(X/S,F)$. Plus gén., le foncteur \deleted{à \unsure{tout} $S'$ \unsure{sur} $(Sch)/S$ qui associe} $S'\longmapsto\Gamma(S',F_{X_{S'}/S'})$ sur $(Sch)/S$ est un faisceau pour fpqc. En effet, on est ramené à le vérifier pour les deux types de familles couvrantes, Zariski, \unsure{et recouvrements|recouvrement} fid plats \deleted{\unsure{quasi-compacts}} \uuline{affines}}

\uline{Corollaire 3} Mettons sur $(Sch)$ une topologie $T$ intermédiaire entre celle de Zariski et la fpqc. Alors pour tt faisceau qu. coh. $F$ sur un $X$, $H^*_{\unsure{1|\,}}(X^T_{\unsure{T|\mathfrak{T}|1}},F^T)\simeq H^*(X,F)$

\uline{Dém} On est ramené par Leray pour $(Sch)^T\longrightarrow(Sch)^{\mathrm{Zar}}$, de prouver que pour \deleted{X} \uline{affine}, \unsure{-- --|on a}.
